\clearpage
\hypertarget{native-468-pallas}{}
\begin{center}
\rule{0.55\linewidth}{0.4pt}\\[2pt]
Astronomische Nachrichten. Band XI. Nr. 262. Seite 403. 9. Juli 1834.\\[2pt]
\rule{0.55\linewidth}{0.4pt}
\end{center}

\begin{center}Opposition der \emph{Pallas} 1834.\end{center}

Von Herrn Hofrat \textsc{Gauss} in Göttingen habe ich die folgenden Beobachtungen der Pallas erhalten, die er mit dem Meridiankreis gemacht hat. Wir haben hier keine erhalten.

\begin{center}
\begin{tabular}{l|l|l}
1834 Febr. 6 & $8^{\rm u}41^{\rm m}47^{\rm s}02$ & $-20^\circ44'23''7$ \\
\hfill 9 & $8\;39\;39.10$ & $-19\;41\;52.9$ \\
\hfill 10 & $8\;38\;57.80$ & $-19\;20\;0.3$ \\
\end{tabular}
\end{center}
\begin{flushright}\textsc{Schumacher.}\end{flushright}

